% =============================================================================
% Section 1: Thesis and Context
% =============================================================================
\section{Thesis and Context}\label{sec:thesis}

\subsection{What the ISI measures}

The International Sovereignty Index (ISI) is a composite indicator that
measures the degree to which a country's bilateral external dependencies
are concentrated among a small number of counterparties.  Higher {\ISI}
scores indicate higher concentration of counterparty exposure and,
consequently, greater structural vulnerability to disruption by any
single partner.

The index aggregates six normalised axes, each constructed as a
Herfindahl--Hirschman Index (\HHI) over the bilateral counterparty
distribution in a specific domain~\cite{isi_paper1}:

\begin{enumerate}[leftmargin=2em]
  \item \textbf{Financial} (cross-border banking and portfolio claims).
  \item \textbf{Energy} (primary energy imports by source country).
  \item \textbf{Technology} (high-technology goods imports).
  \item \textbf{Defence} (arms transfers by supplier).
  \item \textbf{Critical Inputs} (critical raw materials imports).
  \item \textbf{Logistics} (containerised freight and port throughput
        by origin).
\end{enumerate}

\noindent
The composite score is the unweighted arithmetic mean of the six
normalised axis scores:
\begin{equation}\label{eq:composite}
  \text{ISI}_i
  \;=\;
  \frac{1}{K}\sum_{j=1}^{K} A_{ij},
  \qquad K = 6,
\end{equation}
where $A_{ij} \in [0,1]$ is the normalised axis~$j$ score for
country~$i$.  The score is bounded on $[0,1]$; a score of~1
indicates monopoly dependence on a single counterparty across the
axis, while~0 indicates perfect diversification.

Methodology, normalisation, and data sources are documented in
Paper~1~\cite{isi_paper1}.  The full empirical results for all
27~EU member states are reported in Paper~2~\cite{isi_paper2}.

\subsection{What ``variance dominance'' means}

A composite indicator that averages $K$~sub-indicators inherits a
specific variance structure.  The total cross-country variance of the
composite is a function of the variances of the individual axes
\emph{and} all pairwise covariances between them.  When one or two
axes exhibit substantially higher cross-country dispersion than the
rest, those axes mechanically dominate the composite's ability to
differentiate countries from one another.

This paper asks: which axes drive the cross-country dispersion of the
ISI?  The answer determines which structural vulnerabilities actually
distinguish one member state from another---and which are, for ranking
purposes, empirically inert.

\subsection{Why it matters}

Composite indicators are used to inform policy, allocate attention,
and structure public debate.  If the discriminatory power of a
composite rests on a narrow subset of its components, then:

\begin{itemize}[leftmargin=2em]
  \item \textbf{Policy responses} that target the dominant axes will
    have disproportionate ranking effects; reforms targeting
    low-variance axes will not.
  \item \textbf{Measurement quality} depends critically on the data
    integrity of the dominant axes.  If those axes rely on sparse or
    contested data sources, the composite inherits that fragility.
  \item \textbf{Communication} to non-specialist audiences must
    acknowledge that the headline ranking is not an equally weighted
    reflection of all six domains; it is, in practice, a
    defence-and-logistics ranking with a residual contribution from
    the other four.
\end{itemize}

\subsection{Scope and relation to Paper~2}

This paper draws exclusively on the {\ISI} v1.0 data snapshot for the
EU-27, vintage~2024, as reported in Paper~2~\cite{isi_paper2}.  All
empirical anchors---composite values, axis scores, tier
classifications, variance shares, rank correlations, and Monte Carlo
results---are identical to those published in Paper~2.  What differs
is the analytical lens: Paper~2 provides a comprehensive empirical
report; this paper extracts and stress-tests a single structural
finding.
